\documentclass[11pt,a4paper]{article}
\usepackage[margin=21mm]{geometry}
\usepackage{amsmath,amssymb,booktabs,array}
\setlength{\parindent}{0pt}
\setlength{\parskip}{4pt}
\pagestyle{plain}
\begin{document}
\begin{center}{\Large\bfseries Homework 6: Concise Answers}\end{center}

\section*{1 (p.\ 1): Student seniorities}
Assume four independent draws. Let $n_i$ be the number of students from
category $i$, with $\sum_i n_i=4$. The unordered composition has probability
\[
 P(n_1,\ldots,n_5)=\frac{4!}{\prod_i n_i!}\prod_i q_i^{n_i}.
\]
\textbf{Laplace smoothing.} If the given proportions are empirical frequencies
from $N$ students, add-one smoothing gives
\[
 q_i=\frac{Np_i+1}{N+5},\qquad
 (p_1,p_2,p_3,p_4,p_5)=(.24,.22,.20,.18,.16).
\]
The original gives no $N$, so the smoothed numerical probabilities depend
on an unspecified sample size.
\begin{enumerate}
\item[(i)] All four are fourth-year students: $\boxed{q_4^4}$.
\item[(ii)] Two first-year and two second-year students:
$\boxed{\binom42 q_1^2q_2^2=6q_1^2q_2^2}$.
\item[(iii)] The most likely composition is
$\boxed{(n_1,n_2,n_3,n_4,n_5)=(1,1,1,1,0)}$, with probability
$\boxed{24q_1q_2q_3q_4}$.
Indeed, if a category is repeated, there is an unused category. Moving one
student from a repeated category $i$ to an unused category $j$ multiplies
its probability by $n_iq_j/q_i>1$, since $n_i\ge2$ and $q_i/q_j\le1.5$.
Thus the maximum uses four distinct categories, specifically the four largest.
\item[(iv)] The least likely composition is
$\boxed{(0,0,0,0,4)}$, with probability $\boxed{q_5^4}$.
Every product is at least $q_5^4$ and its multinomial coefficient is at least one.
\end{enumerate}
\textbf{Numerical example, assuming $N=100$:}
$q=(25,23,21,19,17)/105$ gives, respectively,
\[
\boxed{0.0010721541,\quad0.0163203603,\quad0.0452996437,\quad0.0006871293.}
\]
For comparison, without smoothing the four probabilities are
\[0.00104976,\quad 0.01672704,\quad 0.04561920,\quad 0.00065536.\]

\section*{2 (p.\ 2): Preparation given a correct answer}
Let $U$ mean unprepared and $R$ mean a correct answer. By Bayes' rule,
\[
 P(U\mid R)=\frac{P(R\mid U)P(U)}{P(R\mid U)P(U)+P(R\mid\neg U)P(\neg U)}
 =\frac{(1/4)(1/4)}{(1/4)(1/4)+(0.9)(3/4)}
 =\boxed{\frac5{59}\approx8.4746\%}.
\]

\newpage
\section*{3 (p.\ 3): Naive Bayes classification}
Use the visible query $x=(\text{yes},\text{yes},\text{yes},\text{yes})$
in the order (Give Birth, Can Fly, Live in Water, Have Legs).
Let $M$ and $N$ denote mammals and non-mammals. There are 7 mammals and
13 non-mammals. The counts of each queried attribute value are:
\begin{center}
\begin{tabular}{c|rrrr|r}
\toprule
Class & Birth=yes & Fly=yes & Water=yes & Legs=yes & Total\\
\midrule
$M$ &6&1&2&5&7\\
$N$ &1&3&3&9&13\\
\bottomrule
\end{tabular}
\end{center}
Using empirical probabilities and conditional independence,
\[
\begin{aligned}
 s_M&=P(M)P(x\mid M)=\frac7{20}\frac67\frac17\frac27\frac57
      =\frac3{343}\approx0.00874636,\\
 s_N&=P(N)P(x\mid N)=\frac{13}{20}\frac1{13}\frac3{13}\frac3{13}\frac9{13}
      =\frac{81}{43940}\approx0.00184342.
\end{aligned}
\]
Therefore
\[
 \boxed{x\text{ is classified as a mammal}},\qquad
 P(M\mid x)=\frac{s_M}{s_M+s_N}=\frac{43940}{53201}\approx0.825924.
\]
If add-one smoothing is also applied here, the binary attributes have two
possible values, while Water has three (yes/no/sometimes). Smoothing the
class priors as well gives
\[
\begin{aligned}
 \widetilde s_M&=\frac8{22}\frac79\frac29\frac3{10}\frac69
     \approx0.01257015,\\
 \widetilde s_N&=\frac{14}{22}\frac2{15}\frac4{15}\frac4{16}\frac{10}{15}
     \approx0.00377104.
\end{aligned}
\]
The smoothed classifier also predicts mammal, with posterior $0.769231$.

\section*{4 (p.\ 4): Weather transition estimates}
Use state order $(S,C,R)$ for rows (current state) and columns (next state).
The 17-day sequence gives 16 adjacent transitions:
\begin{center}
\begin{tabular}{c|rrr|r}
\toprule
From $\backslash$ To & $S$ & $C$ & $R$ & Row total\\
\midrule
$S$&4&1&2&7\\
$C$&0&2&2&4\\
$R$&2&2&1&5\\
\midrule
Total&6&5&5&16\\
\bottomrule
\end{tabular}
\end{center}
Normalize each row to estimate $A_{ij}=P(W_{t+1}=j\mid W_t=i)$:
\[
 \boxed{A=\begin{pmatrix}4/7&1/7&2/7\\0&1/2&1/2\\2/5&2/5&1/5\end{pmatrix}}.
\]
The last day's $C$ has no observed successor and adds no transition.

\newpage
\section*{5 (p.\ 5): Three consecutive sunny days}
The given matrix, in state order $(S,C,R)$, is
\[
 A=\begin{pmatrix}.8&.15&.05\\.38&.6&.02\\.75&.05&.2\end{pmatrix}.
\]
A three-day sunny sequence contains two transitions. For a window starting
on day $t$,
\[
 \boxed{P(W_t=S,W_{t+1}=S,W_{t+2}=S)=P(W_t=S)(.8)^2}.
\]
If day $t$ is already known to be sunny, the conditional probability is $0.64$.
For a window starting on an arbitrary day in the stationary regime, solve
$\pi A=\pi$, $\sum_i\pi_i=1$:
\[
 20\pi_S=38\pi_C+75\pi_R,\qquad
 40\pi_C=15\pi_S+5\pi_R,
 \qquad
 \pi=\frac1{929}(638,245,46).
\]
Hence the stationary answer is
\[
 \boxed{\pi_S(.8)^2=\frac{10208}{23225}\approx0.43952637}.
\]
An initial distribution is not specified. Outside stationarity the general
formula above applies; if $\mu$ is the day-1 distribution, use
$P(W_t=S)=(\mu A^{t-1})_S$.

\section*{6 (p.\ 6): Car-value Bayesian network}
Let $M,E,A,C$ denote Mileage, Engine, Air Conditioner, and Car Value.
The graph gives
\[
 \boxed{P(M,E,A,C)=P(M)P(A)P(E\mid M)P(C\mid E,A)}.
\]
The table represents 40 records. Summing the appropriate counts gives
\[
 P(M=\mathrm{Hi})=P(M=\mathrm{Lo})=\tfrac12,\qquad
 P(A=\mathrm{Working})=\tfrac{25}{40}=\tfrac58,\quad
 P(A=\mathrm{Broken})=\tfrac{15}{40}=\tfrac38.
\]
\begin{center}
\begin{tabular}{c|cc}
\toprule
$M$&$P(E=\mathrm{Good}\mid M)$&$P(E=\mathrm{Bad}\mid M)$\\
\midrule
Hi&$10/20=1/2$&$10/20=1/2$\\
Lo&$15/20=3/4$&$5/20=1/4$\\
\bottomrule
\end{tabular}
\quad
\begin{tabular}{cc|cc}
\toprule
$E$&$A$&$P(C=\mathrm{Hi}\mid E,A)$&$P(C=\mathrm{Lo}\mid E,A)$\\
\midrule
Good&Working&$12/16=3/4$&$4/16=1/4$\\
Good&Broken&$6/9=2/3$&$3/9=1/3$\\
Bad&Working&$2/9$&$7/9$\\
Bad&Broken&$0/6=0$&$6/6=1$\\
\bottomrule
\end{tabular}
\end{center}
For example, marginalizing $M$ gives $P(C=\mathrm{Hi}\mid E=\mathrm{Good},
A=\mathrm{Working})=(3+9)/(3+4+9+0)=3/4$.
Since $E$ and $A$ are marginally independent in this network,
\[
\begin{aligned}
P(E=\mathrm{Bad},A=\mathrm{Broken})
 &=P(A=\mathrm{Broken})\sum_mP(m)P(E=\mathrm{Bad}\mid m)\\
 &=\frac38\left(\frac12\frac12+\frac12\frac14\right)
 =\boxed{\frac9{64}=0.140625}.
\end{aligned}
\]
This is the fitted-network probability; the raw empirical joint frequency
is $(0+4+0+2)/40=0.15$.

\newpage
\section*{7 (p.\ 7): Battery, fuel, gauge and starting}
Let $B,F,G,S$ denote Battery, Fuel, Gauge and Start. The graph has
$B,F$ as independent roots, and both are parents of $G$ and $S$. Therefore
\[
 \boxed{P(B,F,G,S)=P(B)P(F)P(G\mid B,F)P(S\mid B,F)}.
\]
In particular, $G$ and $S$ are conditionally independent given $B,F$.
The supplied root probabilities are
\[
 P(B=\mathrm{bad})=.1,\quad P(B=\mathrm{good})=.9,\qquad
 P(F=\mathrm{empty})=.2,\quad P(F=\mathrm{not\ empty})=.8.
\]
\textbf{(a)} Using the complementary start probability,
$P(S=\mathrm{yes}\mid B=\mathrm{good},F=\mathrm{empty})=1-.8=.2$:
\[
\begin{aligned}
&P(B=\mathrm{good},F=\mathrm{empty},G=\mathrm{empty},S=\mathrm{yes})\\
&\qquad=(.9)(.2)(.8)(.2)=\boxed{0.0288}.
\end{aligned}
\]
\textbf{(b)} Here
$P(G=\mathrm{not\ empty}\mid B=\mathrm{bad},F=\mathrm{empty})=1-.9=.1$:
\[
\begin{aligned}
&P(B=\mathrm{bad},F=\mathrm{empty},G=\mathrm{not\ empty},S=\mathrm{no})\\
&\qquad=(.1)(.2)(.1)(1)=\boxed{0.002}.
\end{aligned}
\]
\textbf{(c)} Marginalize the unknown fuel state. Since $B$ and $F$ are independent,
\[
\begin{aligned}
 P(S=\mathrm{yes}\mid B=\mathrm{bad})
 &=\sum_fP(S=\mathrm{yes}\mid B=\mathrm{bad},f)P(f)\\
 &=(1-.9)(.8)+(1-1)(.2)=\boxed{0.08}.
\end{aligned}
\]
No gauge evidence is given, so its conditional probabilities sum to one
and do not enter this marginal.
\end{document}
